\documentclass[10pt]{article}
\usepackage[a4paper,margin=0.70in]{geometry}
\usepackage{amsmath,amssymb,amsthm,mathtools}
\usepackage{graphicx}
\usepackage[hidelinks]{hyperref}
\setlength{\parindent}{0pt}
\setlength{\parskip}{0.40em}
\newtheorem{theorem}{Theorem}
\newtheorem{corollary}[theorem]{Corollary}
\newtheorem{proposition}[theorem]{Proposition}
\newcommand{\C}{\mathbb C}
\newcommand{\E}{\mathbb E}
\newcommand{\Z}{\mathbb Z}
\newcommand{\HS}{\mathrm{HS}}
\title{A Fourier criterion and phase engineering for\\optimal quantum lifts of abelian Cayley chains}
\author{Candidate substantial partial answer to arXiv:2505.12187}
\date{11 August 2026}

\begin{document}
\maketitle

\textbf{Status.} Candidate substantial partial result, likely valid.  The
open matrix problem is solved for finite abelian Cayley generators; arbitrary
symmetric $Q$-matrices remain open.

\section*{Source question}

Li and Lu \cite{LL} construct a quantum lift of a symmetric finite-state
generator $Q$.  On source PDF page 31, after equations (4.21)--(4.22), they ask
which families of $Q$-matrices admit constants
\[
 M\preceq C_1Q^2+C_2(-Q),\qquad
 \sqrt{C_1}+\sqrt{C_2/\lambda_Q}=O(1),                 \tag{1}
\]
where $M$ is the Gram matrix of $y\mapsto[H,[H,\operatorname{diag}y]]$
and $\lambda_Q$ is the gap of $-Q$.

\begin{center}
\includegraphics[width=0.82\textwidth]{figures/open_problem_crop.png}
\end{center}

\section*{Proof intuition}

The second term in (1) does not enlarge the asymptotic class: on mean-zero
vectors, $\lambda_Q(-Q)\preceq Q^2$.  Thus the problem is a uniform
``discrete Hessian versus Laplacian'' estimate.  For an abelian Cayley chain,
translation invariance diagonalizes this Hessian Gram matrix.  Its Fourier
symbol is unusually simple: an averaged fourth power of differences of the
Hamiltonian symbol.  That identity gives exact constants.

It also exposes a coherent-phase effect.  Positive square roots make all
complete-graph amplitudes add at the trivial character, producing a constant
$|G|/2$.  Cubic phases spread the Fourier mass; Weil's exponential-sum bound
then makes the same complete-graph lift uniformly optimal.  The part still
open is a comparably structural characterization without translation
invariance.

\section*{Exact abelian Fourier criterion}

Let $G$ be a finite abelian group, $N=|G|$, and let
\[
 (Qf)(x)=\sum_{s\in G}q(s)\bigl(f(x+s)-f(x)\bigr),
 \quad q(0)=0,\quad q(-s)=q(s)\geq0,
\]
be ergodic.  Choose $h(0)=0$ with
\[
 h(-s)=\overline{h(s)},\qquad 2|h(s)|^2=q(s),
\]
and let $H_{x,z}=h(z-x)$.  Thus $H=H^*$, and the source's overdamped
calculation gives exactly $Q$; its canonical choice is $h(s)=\sqrt{q(s)/2}$.
For a character $\chi_\eta$, put
\[
 a(\eta)=\sum_{s\in G}h(s)\overline{\chi_\eta(s)},\qquad
 \lambda(\xi)=\sum_sq(s)\bigl(1-\chi_\xi(s)\bigr).
\]
Symmetry makes both quantities real, and $Q\chi_\xi=-\lambda(\xi)\chi_\xi$.
Define $M_h$ by
\[
 \langle f,M_hf\rangle_{\ell^2(G)}
 =\|[H,[H,\operatorname{diag}f]]\|_{\HS,N}^{,2},
\]
using normalized counting and normalized Hilbert--Schmidt norms.

\begin{theorem}[Fourier diagonalization]\label{thm:fourier}
Every character $\chi_\xi$ is an eigenvector of $M_h$, with eigenvalue
\begin{equation}\label{eq:mxi}
 m(\xi)=\frac1N\sum_{\eta\in\widehat G}
             |a(\eta)-a(\eta-\xi)|^4.
\end{equation}
Consequently the least $C$ for which $M_h\preceq C Q^2$ on constants
perpendicular is
\begin{equation}\label{eq:criterion}
 \mathcal C(Q,h)=\max_{\xi\ne0}
 \frac{N^{-1}\sum_\eta|a(\eta)-a(\eta-\xi)|^4}
      {\lambda(\xi)^2}.
\end{equation}
For a family of such generators, the source condition (1) holds uniformly if
and only if $\sup\mathcal C(Q,h)<\infty$.
\end{theorem}

\begin{proof}
Take $f=\chi_\xi$ and write the matrix entry with $z=x+d$ and the intermediate
vertex $x+s$.  A direct expansion gives
\[
 [H,[H,\operatorname{diag}f]]_{x,x+d}=\chi_\xi(x)b_\xi(d),
\]
where
\[
 b_\xi(d)=(1+\chi_\xi(d))(h*h)(d)
             -2\bigl((h\chi_\xi)*h\bigr)(d).
\]
With the displayed convention for $a$, the Fourier transform is the perfect
square
\[
 \widehat b_\xi(\eta)
   =a(\eta)^2+a(\eta-\xi)^2-2a(\eta)a(\eta-\xi)
   =(a(\eta)-a(\eta-\xi))^2.
\]
Translation covariance makes the characters eigenvectors of the Gram
operator.  Parseval now gives \eqref{eq:mxi}, and simultaneous diagonalization
with $Q^2$ gives \eqref{eq:criterion}.

It remains to compare with (1).  Put $L=-Q$.  On constants perpendicular,
$L\succeq\lambda_Q I$, hence $\lambda_QL\preceq L^2$.  If (1) holds and its
displayed square-root sum is at most $B$, then
\[
 M_h\preceq C_1L^2+C_2L\preceq2B^2L^2.
\]
Conversely $M_h\preceq C L^2$ gives (1) with $C_1=C,C_2=0$.
\end{proof}

\section*{Sharp examples}

\begin{corollary}[hypercubes]\label{cor:cube}
For the normalized walk on $G=(\Z/2\Z)^d$, with $q(e_j)=1/d$ and the
canonical positive $h$, one has
\[
                         \mathcal C(Q,h)=3-\frac2d.
\]
Thus this family satisfies the source's optimal-lift criterion uniformly.
\end{corollary}

\begin{proof}
If $\xi$ has Hamming weight $k$, then $\lambda(\xi)=2k/d$ and
\[
 a(\eta)-a(\eta-\xi)=\sqrt{\frac2d}
          \sum_{j\in\operatorname{supp}\xi}(-1)^{\eta_j}.
\]
For a sum $S_k$ of $k$ independent Rademacher variables,
$\E S_k^4=3k^2-2k$.  Hence
\[
 \frac{m(\xi)}{\lambda(\xi)^2}
 =\frac{4(3k^2-2k)/d^2}{4k^2/d^2}=3-\frac2k.
\]
The maximum occurs at $k=d$.
\end{proof}

\begin{proposition}[complete graph: obstruction and repair]\label{prop:complete}
Let $G$ have order $N$ and $q(s)=1/(N-1)$ for every $s\ne0$.

\emph{(i)} For the canonical positive choice
$h(s)=1/\sqrt{2(N-1)}$, one has the sharp identity
\[
                          M_h=\frac N2Q^2.
\]
In particular every pair $C_1,C_2$ in (1) satisfies
$\sqrt{C_1}+\sqrt{C_2/\lambda_Q}\geq\sqrt{N/2}$.

\emph{(ii)} If $G=\Z/p\Z$ with prime $p\geq5$, set, for $s\ne0$,
\[
 h(s)=\frac{\exp(2\pi i s^3/p)}{\sqrt{2(p-1)}}.
\]
Then $H=H^*$, it has the same overdamped complete-graph generator $Q$, and
\[
                              M_h\preceq256Q^2.
\]
Thus Hamiltonian phases restore a dimension-free optimal-lift estimate.
\end{proposition}

\begin{proof}
For (i), the Hamiltonian symbol equals $\sqrt{(N-1)/2}$ at the trivial
character and $-1/\sqrt{2(N-1)}$ elsewhere.  For nontrivial $\xi$, only
$\eta=0,\xi$ contribute to \eqref{eq:mxi}.  Since
$\lambda(\xi)=N/(N-1)$, formula \eqref{eq:criterion} gives exactly $N/2$.
All nonconstant characters have the same eigenvalue, proving the operator
identity and the lower bound in (1).

For (ii), $h(-s)=\overline{h(s)}$ and $2|h(s)|^2=1/(p-1)$.  Its symbol is a
cubic additive-character sum with the $s=0$ term removed.  Weil's bound for a
degree-three polynomial over $\mathbb F_p$ \cite{Weil} yields
\[
 |a(\eta)|\leq\frac{2\sqrt p+1}{\sqrt{2(p-1)}}<2.
\]
Thus \eqref{eq:mxi} gives $m(\xi)<16\cdot2^4=256$.  Here
$\lambda(\xi)=p/(p-1)>1$ for $\xi\ne0$, so
$m(\xi)/\lambda(\xi)^2<256$, and Theorem~\ref{thm:fourier} finishes the proof.
\end{proof}

\newpage
\section*{Verification, limitations, and novelty}

The included script checks the exact complete-graph and hypercube constants
and evaluates the cubic construction for primes through $101$; all assertions
pass.  These computations are regression tests, not proof.  The argument above
uses only Fourier inversion, a Rademacher fourth moment, and the classical
cubic Weil bound.

The result characterizes all abelian Cayley instances but not arbitrary
symmetric matrices.  It also shows that the answer depends on Hamiltonian
phases: part (i) concerns exactly the positive square root in source equation
(4.17), while part (ii) uses the more general self-adjoint Hamiltonian allowed
by the source's preceding overdamped construction.

\subsection*{Consequences for the source problem}

First, within the abelian Cayley class the two constants in (4.21) carry no
more asymptotic information than the single fourth-difference ratio
\eqref{eq:criterion}; the $C_2$ term can always be absorbed into $C_1Q^2$ at a
universal loss once (4.22) is imposed.  Thus \eqref{eq:criterion} is an exact
and directly computable answer for this class, not merely a sufficient test.

Second, small spectral gap and large degree are neither individual
obstructions.  Hypercubes have gap $2/d$ and degree $d$ but a constant below
$3$, whereas the normalized complete graph has gap tending to one but the
canonical constant is $N/2$.  The relevant datum is coherent concentration of
the square-root Hamiltonian symbol.

Third, a characterization purely in terms of $Q$ necessarily presupposes a
phase convention such as source equation (4.17).  The same complete-graph
$Q$ fails for the positive convention and succeeds for cubic phases.  This
suggests separating two future questions: classify the canonical lift, and
optimize over all self-adjoint Hamiltonians with $2|h_{xz}|^2=Q_{xz}$.

On 11 August 2026, the run registry, solution, attempt, and proof-gap indexes
were searched for the arXiv id and core quantum-lift/Fourier terms; no answer
was found.  No broad external literature search was performed, so novelty
confidence is moderate-to-low.

\textbf{Human-review recommendation.} Check the Fourier shift convention in
\eqref{eq:mxi}, the deletion of the zero term in the Weil estimate, and the
application of the source's abstract lift theorem to the phase-modified
self-adjoint Hamiltonian.  The general non-Cayley problem remains open.

\begin{thebibliography}{9}
\bibitem{LL}
B. Li and J. Lu,
\emph{Speeding up quantum Markov processes through lifting},
arXiv:2505.12187 (2025), especially Section 4.2 and equations (4.17)--(4.22).

\bibitem{Weil}
A. Weil,
\emph{On some exponential sums},
Proc. Natl. Acad. Sci. USA 34 (1948), 204--207.
\end{thebibliography}

\end{document}
